
\documentclass[12pt]{article}
\usepackage{amsmath}
\usepackage{listings}
\usepackage{xcolor}
\usepackage{geometry}
\geometry{margin=1in}

\title{Advanced Python-Based Gene Sequencing Pipeline\\Detailed Syntax and Workflow Explanation}
\author{}
\date{}

\lstset{
    basicstyle=\ttfamily\small,
    keywordstyle=\color{blue},
    commentstyle=\color{gray},
    stringstyle=\color{red},
    breaklines=true,
    frame=single
}

\begin{document}

\maketitle

\section{Introduction}
This document provides a detailed breakdown of a Python function used in sequencing data preprocessing, followed by an extended explanation of a full computational gene sequencing workflow.

\section{Core Function}
\begin{lstlisting}[language=Python]
def trim_read(seq, qual, threshold=20):
    return "".join([base for base, q in zip(seq, qual) if q >= threshold])
\end{lstlisting}

\section{Step-by-Step Syntax Explanation}

\subsection{Function Definition}
\begin{itemize}
\item \texttt{def}: घोषित करता है कि हम एक function define कर रहे हैं।
\item \texttt{trim\_read}: function का नाम।
\item \texttt{seq}: DNA sequence (string).
\item \texttt{qual}: quality scores (list of integers).
\item \texttt{threshold=20}: default parameter।
\end{itemize}

\subsection{zip Function}
\begin{lstlisting}[language=Python]
zip(seq, qual)
\end{lstlisting}
Pairs corresponding elements:
\begin{verbatim}
"ATCG" + [30,10,25,18] → [('A',30), ('T',10), ...]
\end{verbatim}

\subsection{List Comprehension}
\begin{lstlisting}[language=Python]
[base for base, q in zip(seq, qual) if q >= threshold]
\end{lstlisting}

Equivalent to:
\begin{lstlisting}[language=Python]
result = []
for base, q in zip(seq, qual):
    if q >= threshold:
        result.append(base)
\end{lstlisting}

\subsection{Join Operation}
\begin{lstlisting}[language=Python]
"".join(list)
\end{lstlisting}
Concatenates characters into a string efficiently.

\section{Execution Example}
\begin{lstlisting}[language=Python]
seq = "ATCG"
qual = [30, 10, 25, 18]
\end{lstlisting}

Output:
\begin{verbatim}
"AC"
\end{verbatim}

\section{Advanced Optimization}
Use generator expression:
\begin{lstlisting}[language=Python]
return "".join(base for base, q in zip(seq, qual) if q >= threshold)
\end{lstlisting}

\section{Full Sequencing Workflow (Computational)}

\subsection{1. FASTQ Parsing}
\begin{lstlisting}[language=Python]
from Bio import SeqIO
for record in SeqIO.parse("reads.fastq", "fastq"):
    seq = record.seq
\end{lstlisting}

\subsection{2. Quality Control}
Filtering reads based on Phred score:
\begin{equation}
Q = -10 \log_{10}(P)
\end{equation}

\subsection{3. Alignment}
\begin{lstlisting}[language=Python]
def naive_align(read, ref):
    for i in range(len(ref)-len(read)+1):
        if ref[i:i+len(read)] == read:
            return i
\end{lstlisting}

\subsection{4. BAM Processing}
\begin{lstlisting}[language=Python]
import pysam
bam = pysam.AlignmentFile("file.bam", "rb")
\end{lstlisting}

\subsection{5. Variant Calling}
\begin{lstlisting}[language=Python]
for pileupcolumn in bam.pileup():
    print(pileupcolumn.pos)
\end{lstlisting}

\subsection{6. Filtering Variants}
\begin{lstlisting}[language=Python]
if quality > 30:
    keep = True
\end{lstlisting}

\subsection{7. Annotation}
\begin{lstlisting}[language=Python]
import vcf
\end{lstlisting}

\subsection{8. Assembly (De Bruijn Graph)}
\begin{lstlisting}[language=Python]
def build_graph(reads, k):
    pass
\end{lstlisting}

\subsection{9. Visualization}
\begin{lstlisting}[language=Python]
import matplotlib.pyplot as plt
plt.plot([1,2,3])
\end{lstlisting}

\subsection{10. Workflow Automation}
\begin{lstlisting}[language=Python]
rule align:
    shell: "bwa mem ref.fa reads.fastq"
\end{lstlisting}

\section{Conclusion}
Python acts as a control layer integrating multiple bioinformatics tools, enabling scalable sequencing analysis.

\end{document}
